\documentclass{article}
\usepackage{amsmath}
\usepackage{geometry}
\geometry{margin=1in}
\begin{document}

\subsection*{Heterogeneous Deterrence Effects of Unemployment and Death Risk}

\paragraph{Unemployment deters the poor (but not the rich).}
Although a proportional income loss due to unemployment generates an identical drop in log consumption utility for all individuals (since $\log(\lambda y) = \log(\lambda) + \log(y)$), its impact is asymmetric across the income distribution due to budget constraints. Low-income individuals are more likely to be close to violating the feasibility constraint $c = y - p_v v > 0$. When income falls, these individuals may no longer be able to afford drug consumption without inducing negative or minimal consumption utility. In contrast, high-income individuals can absorb the same proportional income loss while maintaining feasible and utility-generating levels of both $c$ and $v$. As a result, drug use becomes unattractive for the poor under unemployment, while the rich are relatively unaffected. Thus, unemployment selectively deters drug use among low-income individuals.

\vspace{1em}

\paragraph{Death risk deters the rich when $b$ is low, and the poor when $b$ is high.}
When drug use carries a probability $\pi$ of death, expected utility becomes:
\[
U_{\text{drug}} = (1 - \pi)\left[\theta \log(c) + (1 - \theta)\log(v + b) + D\right]
\]
where $D > 0$ captures the value of being alive (e.g., from family, purpose, or general life satisfaction). Upon death, utility falls to zero. By contrast, sober living guarantees survival and yields:
\[
U_{\text{sober}} = \theta \log(y) + (1 - \theta)\log(b) + D
\]

When $b$ is low, sober utility is very low for the poor due to small $y$ and low $b$. Drug use, even with a chance of death, may yield higher expected utility, especially since the value-of-life term $D$ is small relative to the utility gain from $\log(v + b)$. Rich individuals, however, enjoy high sober utility and have more to lose from the potential loss of $D$, making them more sensitive to death risk.

When $b$ is high, even poor individuals derive substantial utility from sober living through $\log(b)$, making the loss of $D$ under death more significant relative to the gains from drug use. In this case, the death penalty effectively reduces the expected utility from drug use below that of sober living, deterring the poor. Rich individuals, in contrast, face high utility in both scenarios. However, because $c$ and $v + b$ both increase with income, the utility from drug use also grows with $y$, and the fixed cost of losing $D$ becomes relatively small. As a result, the rich may still prefer drug use despite the death risk, as the penalty constitutes a smaller share of their total utility.


\end{document}

